\section{\personMaslov{}--\personMorse{}指标}
\label{sec:maslov-morse-index}

\stfhIndexRangeStart{I127}{Masiluofu-Moersizhibiao}{马斯洛夫--莫尔斯指标（Maslov--Morse index）}
本节要给 $\mathcal A$ 的每个临界点赋予一个\termAbsoluteMorseIndex{}。一般而言，
赋予这样的绝对指标是一件微妙的事，因为 $\mathcal A$ 在任意临界点处通常的\personMorse{}指标都是无穷。我们把这里所要定义的\personMorse{}指标称为该临界点的\termMaslovMorseIndex{}。这个指标的定义与 $\mathcal A$ 的定义有几分相似。不过，
为定义它，还需要固定 $\ell_0^*\Lambda(M)$ 的一个截面 $\lambda^0$，使得
\[
  \lambda^0(0)=T_{\ell_0(0)}L_0,
  \qquad
  \lambda^0(1)=T_{\ell_0(1)}L_1.
\]
这里，$\Lambda(M)$ 表示 $TM$ 的拉格朗日--\personGrassmann{}流形丛：
\[
  \Lambda(M)=\bigcup_{x\in M}\Lambda(T_xM),
\]
其中
\[
  \Lambda(T_xM):=\mathcal L(T_xM,\omega_x)
\]
是辛向量空间 $(T_xM,\omega_x)$ 的拉格朗日子空间集合。（这里，为避免与下一节
的拉格朗日链 $\mathcal L$ 在记号上混淆，我们更换了拉格朗日--\personGrassmann{}流形的
记号。）

对临界点 $[\widehat x,w]$，我们要在 $[0,1]^2$ 上关联一个连续辛丛对
$(E_w,\lambda_w)$。利用第一卷命题 3.2.14，可像前面一样构造辛丛对
$(E_w,\lambda_w)$；它在同伦意义下唯一确定。先取
\[
  E_w=w^*TM.
\]
为定义 $\lambda_w$，选取沿
\[
  \Lambda_1(T_xM,\omega_x;T_xL_0)
  \subset\Lambda(T_xM,\omega_x)
\]
的一个向量场 $\xi^x$，使得（见第一卷引理 3.2.12）
\[
  \xi^x(T_xL_0)\in C_0(T_xL_0).
\]
\footnote{译者注：原书在上式中写成 $\xi^x\in C_0(T_xL_0)$；但 $\xi^x$
是沿马斯洛夫闭链的向量场，而 $C_0(T_xL_0)$ 是拉格朗日--\personGrassmann{}流形在
$T_xL_0$ 处切空间中的锥。下一行也明确使用 $\xi^x(T_xL_0)$，译文依此补上取值。}
再选取一条满足第一卷命题 3.2.14 条件的道路
\[
  \alpha^x:[0,1]\longrightarrow\Lambda(T_xM,\omega_x),
\]
特别地，
\[
  \alpha^x(0)=T_xL_0,
  \qquad
  \alpha^x(1)=T_xL_1\subset T_xM,
\]
并且
\[
  (\alpha^x)'(0)=\xi^x(T_xL_0).
\]
这里还要求
\[
  \alpha^x(t)\in
  \Lambda_0(T_xM,\omega_x;T_xL_0)
  :=\Lambda(T_xM,\omega_x)
    \setminus\Lambda_1(T_xM,\omega_x;T_xL_0),
  \qquad 0<t\leq1.
\]
\footnote{译者注：原书把 $\alpha^x$ 的靶写成未在本节引入的
$\Lambda(S,\Omega)$；当前辛向量空间是 $(T_xM,\omega_x)$，译文据此补正。
此外，原书虽引用第一卷命题 3.2.14 并随即使用其同伦唯一性，却只明写了端点和
初始切向量；该命题还要求道路在 $0<t\leq1$ 时避开相对于 $T_xL_0$ 的马斯洛夫
闭链。译文将这项必要条件明示。}
该命题还说明，$(\xi^x,\alpha^x)$ 的任意两个这种选取彼此同伦。

现在考虑 $E_w|_{\partial[0,1]^2}$ 的连续拉格朗日子丛
\[
  \lambda_w\longrightarrow\partial[0,1]^2,
\]
它在 $\partial[0,1]^2$ 每一点处的纤维由下式给出：
\[
  \begin{aligned}
    \lambda_w(s,0)&=T_{w(s,0)}L_0,
    &\lambda_w(1,t)&=\alpha^x(t),\\
    \lambda_w(s,1)&=T_{w(s,1)}L_1,
    &\lambda_w(0,t)&=\lambda^0(t).
  \end{aligned}
\]
\footnote{译者注：原书在这里把此前唯一引入的 $E_w$ 写成未定义的 $E$，又把
区间 $[0,1]$ 上的截面值 $\lambda^0(t)$ 写成 $\lambda^0_{(0,t)}$。译文依丛及
截面的定义域补正。}
第一卷命题 3.2.14 再次蕴含，$\lambda_w$ 在同伦意义下唯一确定。

\begin{definition}
\label{def:maslov-morse-index}
定义 $[\widehat x,w]$ 的\termMaslovMorseIndex{}为
\[
  \mu\bigl([\widehat x,w];(\ell_0,\lambda^0)\bigr)
  =\mu(E_w,\lambda_w).
\]
\end{definition}

注意，这个绝对指标依赖于 $(\ell_0,\lambda^0)$ 以及
$[w]\in\pi_2(x;\ell_0)$ 的选取。\footnote{译者注：原书在这里把已固定的
$\lambda^0$ 改写成 $\lambda_0$，并把上一节定义的 $\pi_2(x;\ell_0)$ 简写成未另行
定义的 $\pi_2(x)$；命题 13.6.3 及其证明又继续混用 $\lambda_0$ 与 $\lambda^0$。
译文统一沿用定义 13.6.1 的 $\lambda^0$ 和上一节的 $\pi_2(x;\ell_0)$。}
\stfhIndex{I128}{Masiluofu-Weitaibozhibiao}{马斯洛夫--维泰博指标（Maslov--Viterbo index）}
下面这个指标的相对版本称为\termMaslovViterboIndex{}；它不需要任何这类选取，
而只依赖于同伦类 $B\in\pi_2(x,y)$。设
\[
  u:[0,1]^2\longrightarrow M
\]
是代表同伦类 $B$ 的光滑映射。

选取 $u^*TM$ 的一个平凡化，并考虑丛对 $(E_u,\lambda_u)$，其中
$\lambda_u$ 是由下式定义的拉格朗日子丛：
\[
  \begin{aligned}
    \lambda_u(s,0)&=T_{u(s,0)}L_0,
    &\lambda_u(1,t)&=\alpha^y(t),\\
    \lambda_u(s,1)&=T_{u(s,1)}L_1,
    &\lambda_u(0,t)&=\alpha^x(t).
  \end{aligned}
\]

\begin{definition}[\personMaslov{}--\personViterbo{}指标]
\label{def:maslov-viterbo-index}
设 $B\in\pi_2(x,y)$，并设 $u$ 如上。以 $\mu(L_0,L_1;B)$ 表示的\personMaslov{}--\personViterbo{}指标，定义为丛对 $(E_u,\lambda_u)$ 的马斯洛夫指标。
\end{definition}

需要指出，这个相对指标的定义不需要作任何选取。指标
$\mu(E_u,\lambda_u)$ 的这个定义恰好就是 (Vi88) 中\personViterbo{}的定义。

下面的命题把\personMaslov{}--\personMorse{}指标与\personMaslov{}--\personViterbo{}指标联系起来；它直接由
定义得到。

\begin{proposition}
\label{prop:maslov-morse-viterbo-difference}
设
\[
  \widehat x,\widehat y\in\Omega(L_0,L_1;\ell_0),
  \qquad
  B\in\pi_2(x,y).
\]
则
\[
  \mu(L_0,L_1;B)
  =\mu\bigl([\widehat y,w\#B];(\ell_0,\lambda^0)\bigr)
   -\mu\bigl([\widehat x,w];(\ell_0,\lambda^0)\bigr).
\]
特别地，右端的差不依赖于 $(\ell_0,\lambda^0)$ 或 $w$ 的选取。
\end{proposition}

\begin{proof}
选取 $u^*TM$ 的一个平凡化，并考虑上面定义的丛对 $(E_u,\lambda_u)$。
用按上述方式构造的丛对 $(E_w,\lambda_w)$，把
\[
  \mu\bigl([\widehat x,w];(\ell_0,\lambda^0)\bigr)
\]
表示为 $\mu(E_w,\lambda_w)$。再把
\[
  \mu\bigl([\widehat y,w\#B];(\ell_0,\lambda^0)\bigr)
\]
表示为 $\mu(E_{w\#u},\lambda_{w\#u})$；这里
$(E_{w\#u},\lambda_{w\#u})$ 是把上述两个丛对拼接起来所得的、在同伦意义下唯一
的丛对。由构造，它们满足
\[
  (E_{w\#u},\lambda_{w\#u})
  =(E_w,\lambda_w)\#(E_u,\lambda_u).
\]
由于马斯洛夫指标对丛对的拼接可加，故
\[
  \mu(E_{w\#u},\lambda_{w\#u})
  =\mu(E_w,\lambda_w)+\mu(E_u,\lambda_u),
\]
或者等价地，
\[
  \mu(E_u,\lambda_u)
  =\mu(E_{w\#u},\lambda_{w\#u})-\mu(E_w,\lambda_w).
\]
命题得证。
\end{proof}
\stfhIndexRangeEnd{I127}{Masiluofu-Moersizhibiao}{马斯洛夫--莫尔斯指标（Maslov--Morse index）}
